\ifintegriereaufgaben
\section*{Integrationsregeln Aufgaben}\label{MA2-integrationsregeln-aufgaben}

Erkl�rungen und Beispiele, die zur Bearbeitung der Aufgaben hilfreich sind, finden Sie im Skript\wennreferenzvorhanden{MA2-integrationsregeln}{ (Abschnitt \abschnittsverweis{MA2-integrationsregeln})} und, falls Sie nicht weiterkommen, schauen Sie \hyperref[MA2-integrationsregeln-loesungen]{hier}.
\else
\section{Integrationsregeln Aufgaben}\label{MA2-integrationsregeln-aufgaben}

Regeln und Beispiele zu diesen Aufgaben finden Sie im Skript. Die L�sungen finden Sie \hyperref[MA2-integrationsregeln-loesungen]{hier}.
\fi

\begin{exercisebox}[St�ckweise stetige Funktionen]
	Gegeben seien die folgenden st�ckweise definierten Funktionen
	\begin{eqnarray*}
		f:x\mapsto \begin{cases} x, \quad &x<1\,,\\ x-1, \quad &x\geq 1 \end{cases}\quad \quad g:x\mapsto \begin{cases} x, \quad & x<1\,,\\ 2, \quad & x=1\,,\\ \cos(x),\quad & x>1 \end{cases}
	\end{eqnarray*}
			\begin{itemize}
				\item[(a)] Integrieren Sie diese Funktionen st�ckweise.
				\item[(b)] Sind die Funktionen, die Sie unter \( (a)\) entwickelt haben Stammfunktionen?
			\end{itemize}
\end{exercisebox}

\begin{exercisebox}[Partielle Integration] % Papula 1, V, Abschnitt 8, Aufgabe %/6
	Berechnen Sie die folgenden Integrale:
	\begin{eqnarray*}i
		(a)\; \int x \ln(x) \, \mathrm{d}x \quad
		(b)\; \int_1^5 \ln(t) \, \mathrm{d}t \quad
		(c)\; \int_0^{0,8} x e^x \, \mathrm{d}x\quad
		(d)\; \int e^x \cos(x) \, \mathrm{d}x
	\end{eqnarray*}
	Finden Sie eine Rekursionsformel zur Berechnung von \[ \int x^n\cdot e^x \, \mathrm d x, \quad n\in \mathbb N^+\,.\]
\end{exercisebox}

\begin{exercisebox}[Substitution] % Papula 1, V Abschnitt 8, Aufg. 1
	Berechnen Sie die folgenden Integrale:
	\begin{eqnarray*}
		(a)\; \int x \sin(x^2) \, \mathrm{d}x \quad
		(b)\; \int_{-1}^{1} \frac{t}{\sqrt{1+t^2}}\, \mathrm{d}t \quad
			(c)\; \int_{-1}^{1} \frac{5+x}{5-x} \, \mathrm{d}x\quad
			(d)\; \int \frac{\sqrt{4-x^2}}{x^2} \, \mathrm{d}x
		\end{eqnarray*}
\end{exercisebox}

\begin{exercisebox}[Integrationsregeln] % Papula 1, V Abschnitt 8, Aufg. 10
	Berechnen Sie die folgenden Integrale:
	\begin{eqnarray*}
		\begin{array}{rl rl rl}
			(a)\; &\displaystyle{\int} \mathrm{cot}(x) \, \mathrm{d}x \quad
			& (b)\;&\displaystyle{\int} x\mathrm{cosh}(x) \, \mathrm{d}x \quad & (c) \; & \displaystyle{\int} \displaystyle{\frac{(\ln(x))^3}{x}}\,\mathrm{d}x \\
			(d)\;& \displaystyle{\int} \displaystyle{\frac{12 x^2}{2x^3-1}} \, \mathrm{d}x  &(e)\; &\displaystyle{\int} \sqrt{x^2 - 2x} \, \mathrm{d}x \quad & (f)\;&\displaystyle{\int} \mathrm{arctan}(x) \, \mathrm{d}x
		\end{array}
	\end{eqnarray*}
\end{exercisebox}

\begin{exercisebox}[Integrationsregeln]
	Berechnen Sie die folgenden Integrale:
	\begin{eqnarray*}
		\begin{array}{rl rl}
			(a)\; &\displaystyle{\int_{1/2}^1} \displaystyle{\frac{\sqrt{\frac1x + 1}}{x^2}} \, \mathrm{d}x, \quad y=\displaystyle{\frac1x} +1
			& (b)\;&\displaystyle{\int_0^{1/b}} \displaystyle{\frac{ax}{bx+c}} \, \mathrm{d}x,  \quad b\not=0, c>0, x = \displaystyle{\frac{t-c}{b}} \\
				(c) \; & \displaystyle{\int} \displaystyle{\frac{1}{\sqrt{x^2-4}}}\,\mathrm{d}x &(d)\;& \displaystyle{\int} \displaystyle{\sqrt{18 + 2x^2}} \, \mathrm{d}x
		\end{array}
	\end{eqnarray*}
\end{exercisebox}

%\begin{exercisebox}[Gebrochenrationale Funktionen]
%	Berechnen Sie folgende Integrale. Hinweis: Nutzen Sie die Ergebnsise auf Blatt $1$.
%	\begin{eqnarray*}
%		\begin{array}{l l l}
%			(a)\; \displaystyle{ \int \frac{-x^4-2x^3+7x^2-5x+3}{ x^2+3x-4}} \, \mathrm{d}x \quad & (b)\; \displaystyle{\int \frac{x^2+x+1}{x^3-3x+2}}\,\mathrm{d}x  & (c) \; \displaystyle{\int \frac{1}{x^3+x}}\,\mathrm{d}x \\[1ex]
%			(d)\; \displaystyle{\int \frac{x^2-5x}{(x-2)^3}} \,\mathrm{d}x &(e)\; \displaystyle{ \int \frac{1}{ x^2-a^2}} \, \mathrm{d}x, \; a>0 & (f)\; \displaystyle{\int \frac{4x^3}{x^3+2x^2-x-2}}\,\mathrm{d}x \\[2ex]
%			(g) \; \displaystyle{\int \frac{3z}{z^3-3z^2-4}}\,\mathrm{d}x  & (h)\; \displaystyle{\int \frac{4x-2}{x^2-2x-63}}\, \mathrm{d}x \quad & (i)\; \displaystyle{\int \frac{2x+1}{x^3-6x^2+9x}}  \, \mathrm{d}x
%		\end{array}
%	\end{eqnarray*}
%\end{exercisebox}

\begin{exercisebox}{Kurvendiskussion mit Fl�chenberechnung}{ } %p. 176
		Gegeben ist die Funktion \(f\) mit \(f:x\mapsto x^2(x^2-4)\). Ihr Graph sei \(K\).
		\begin{itemize}
			\item[(a)] Untersuchen Sie \(f\) und zeichenn Sie den Graphen.
			\item[(b)] Wie gro� ist der Inhalt der Fl�che, die \(K\) mit der \(x\)-Achse einschlie�t?
			\item[(c)] Dei Tangente in den Tiefpunkten von \(K\) begrenzt mit \(K\) eine Fl�che. Berechnen Sie den Inhalt dieser Fl�che.
		\end{itemize}
\end{exercisebox}



